\section{Short-pulse frequency inference and calibration}
\label{sec:inference-calibration}

We now specify the local \textsc{Track} branch introduced in Sec.~II.C.
A usable tracking iteration must convert the differential population into a
frequency estimate, move the physical working point toward
$\omega_{\mathrm{tar}}$, and keep the next measurement within the validated
response interval as the qubit frequency changes. We therefore derive the
short-pulse differential response, construct its nonlinear local inverse
together with an explicit validity boundary, use the resulting frequency
estimate in a bounded flux-bias update, and update the drive reference to
preserve the local measurement condition. Together, these elements specify how, and within what limits, short-pulse
frequency information can be incorporated into closed-loop flux-bias
calibration.


\subsection{Short-pulse differential response}
\label{sec:differential-response}

To characterize how detuning is encoded in
Eq.~\eqref{eq:pdiff}, we regard the differential population as a
functional of a time-dependent detuning $\Delta(t)$. Following the general
framework of linear and nonlinear response
theory~\cite{kubo1957statistical,mukamel1995principles}, we expand the
differential population around the resonant control evolution as
%
\begin{equation}
\begin{aligned}
p_{\mathrm d}[\Delta]
={}&
p_{\mathrm d}^{(0)}
+
\sum_{m=1}^{\infty}
\frac{1}{m!}
\int_0^T\!\dd t_1\cdots\int_0^T\!\dd t_m
\\
&\quad\times
k_m(t_1,\ldots,t_m)
\prod_{j=1}^{m}\Delta(t_j).
\end{aligned}
\label{eq:differential-volterra}
\end{equation}
%
Here $k_m(t_1,\ldots,t_m)$ is the symmetrized $m$th-order
detuning-response kernel of the complete bipartite short-pulse probe. It
quantifies the joint contribution of detuning at the times
$t_1,\ldots,t_m$ to the final differential population, and $T$ is the total
sequence duration. The response-function derivation, including time ordering
and nested commutators, is given in the Supplemental Material.

For scalar local frequency estimation, we restrict attention to the
quasi-static regime in which detuning variations during a single short
sequence are negligible, so that $\Delta(t)\simeq\Delta$. This approximation
does not preclude drift between repetitions or feedback iterations.

For the ideal $+X/-X$ probe defined in
Eq.~\eqref{eq:bipartite-sequences}, detuning reversal interchanges the branch
responses,
$p_{+X}(-\Delta)=p_{-X}(\Delta)$. The differential combination therefore
satisfies $p_{\mathrm d}(-\Delta)=-p_{\mathrm d}(\Delta)$. In particular,
$p_{\mathrm d}^{(0)}=0$, and all even-order terms in
Eq.~\eqref{eq:differential-volterra} vanish. The local response can thus be
written as
%
\begin{equation}
p_{\mathrm d}(\Delta)
=
G_1\Delta
+
\frac{G_3}{3!}\Delta^3
+
\mathcal O(\Delta^5),
\label{eq:short-pulse-cubic-response}
\end{equation}
%
where $G_1$ and $G_3$ are the first and third derivatives of
$p_{\mathrm d}$ with respect to $\Delta$ at resonance. The coefficient
$G_1$ determines the local sensitivity, while $G_3$ gives the leading
nonlinear correction.

For the experimental implementation, we separate the response measurements
into calibration and validation data. From the response-calibration data, we
obtain $G_1$ and $G_3$ through an offline, measurement-based kernel
calibration. Calibrated Virtual-$Z$ phase insertions are applied at selected
times within the two control branches, and finite differences of the
resulting branch populations sample the detuning-response derivatives. In
the quasi-static limit, the integrated coefficients are
%
\begin{equation}
\begin{aligned}
G_1
&=
\int_0^T k_1(t)\,\dd t,
\\
G_3
&=
\iiint_{[0,T]^3}
k_3(t_1,t_2,t_3)\,
\dd t_1\dd t_2\dd t_3 .
\end{aligned}
\label{eq:g1-g3}
\end{equation}
%
The Virtual-$Z$ operations follow
Ref.~\cite{mckay2017efficient}. The finite-difference construction and
three-time sampling scheme used to obtain the response kernels are detailed
in the Supplemental Material.

Experimental imperfections may nevertheless produce a nonzero measured
differential signal at resonance. We therefore define
%
\begin{equation}
\widetilde p_{\mathrm d}
=
p_{\mathrm d}^{\mathrm{meas}}-b_0,
\label{eq:corrected-differential-response}
\end{equation}
%
where $b_0$ is the zero-detuning intercept estimated from the
response-calibration data and held fixed thereafter. This operation corrects
a differential-population offset; it does not shift the frequency reference.

For comparison, an odd-polynomial fit to the corrected quasi-static response
provides a separate check of the integrated coefficients but is not used to
generate the coefficients supplied to the main feedback protocol. Once
$G_1$, $G_3$, and $b_0$ have been fixed, separately acquired
response-validation data are used to establish
$\Delta_{\mathrm{val}}$ and to assess residual symmetry breaking through
$p_{\mathrm{even}}(\Delta)
=[\widetilde p_{\mathrm d}(\Delta)
+\widetilde p_{\mathrm d}(-\Delta)]/2$.
The even component is treated as a diagnostic of model mismatch rather than
absorbed into the ideal odd response model. This response validation is
distinct from the independent frequency verification used to confirm
calibration success in Sec.~II.C.

The acquisition and partitioning of the calibration and validation data,
together with the sampling schedule, coefficient-reuse range, recalibration
conditions, and associated cost accounting, are specified in
Appendix~\ref{app:methods}.

\subsection{Local frequency estimator}
\label{sec:local-estimator}

The first-order frequency estimate follows from the local slope:
%
\begin{equation}
\widehat\Delta^{(1)}
=
\frac{\widetilde p_{\mathrm d}}{G_1},
\qquad
\widehat\omega_{01}^{(1)}
=
\omega_d+\widehat\Delta^{(1)}.
\label{eq:linear-frequency-estimator}
\end{equation}
%
Its truncation bias increases as the probe moves away from resonance. This
estimate provides the initial value for the cubic inversion and the
first-order reference used to quantify nonlinear-estimation bias in
Sec.~IV.

The main local estimator retains the third-order response and solves
%
\begin{equation}
G_1\widehat\Delta^{(3)}
+
\frac{G_3}{6}
\left(\widehat\Delta^{(3)}\right)^3
=
\widetilde p_{\mathrm d},
\label{eq:cubic-frequency-estimator}
\end{equation}
%
with
$\widehat\omega_{01}^{(3)}
=\omega_d+\widehat\Delta^{(3)}$.
Newton iteration is initialized with $\widehat\Delta^{(1)}$, and the
candidate solution is selected from the branch connected continuously to
$\Delta=0$.

When $G_1G_3<0$, the cubic response has stationary points at
%
\begin{equation}
\left|\Delta_{\mathrm{fold}}\right|
=
\sqrt{-\frac{2G_1}{G_3}}.
\label{eq:response-fold}
\end{equation}
%
Beyond these points, the cubic response is no longer one-to-one, and the same
differential population can correspond to multiple detunings. The response
fold is therefore a mathematical boundary of the truncated cubic model, not
an experimentally established operating range.

We define the validated local operating interval by
%
\begin{equation}
|\Delta|
\leq
\Delta_{\mathrm{val}},
\qquad
\Delta_{\mathrm{val}}
<
|\Delta_{\mathrm{fold}}|
\quad
(G_1G_3<0).
\label{eq:validated-interval}
\end{equation}
%
The value of $\Delta_{\mathrm{val}}$ is established from the held-out
response-validation data using predefined estimator-error and confidence
criteria. It remains necessary when the cubic model has no real fold because
model truncation, measurement noise, and experimental imperfections can
restrict the usable response before a mathematical turnover occurs. Failure
of numerical convergence, branch continuity, or the validated-range check
rejects the estimate and triggers \textsc{Reacquire} rather than extrapolation
of the local inverse.

\subsection{Track initialization and flux-bias update}
\label{sec:calibration-protocol}

Section~II.C defined the state transitions of the complete calibration
workflow. Here we construct the local controller that operates after an
\textsc{Acquire} or \textsc{Reacquire} measurement has supplied the seed
%
\begin{equation}
\left(
\Phi_{\mathrm{seed}},
\widehat\omega_{01,\mathrm{seed}},
\sigma_{\omega,\mathrm{seed}}
\right).
\label{eq:acquired-seed}
\end{equation}
%
The acquired frequency estimate is reused rather than repeated with the
short-pulse probe at the same bias. It initializes the local drive reference
as
$\omega_{d,\mathrm{seed}}
=\widehat\omega_{01,\mathrm{seed}}$ and gives the seed residual
%
\begin{equation}
\widehat r_{\mathrm{seed}}
=
\widehat\omega_{01,\mathrm{seed}}
-
\omega_{\mathrm{tar}}.
\label{eq:seed-residual}
\end{equation}
%
If Eq.~\eqref{eq:target-entry-condition} is already satisfied, the workflow
proceeds directly to \textsc{Lock}. Otherwise, the seed is used to initialize
the local flux-bias update.

Because one bias--frequency pair does not determine a secant sensitivity, the
controller proposes a bounded one-sided step,
%
\begin{equation}
\Phi_1^{\mathrm{prop}}
=
\Phi_{\mathrm{seed}}
-
\chi_\Phi
\operatorname{sgn}(\widehat r_{\mathrm{seed}})
\delta\Phi_{\mathrm{blind}},
\label{eq:blind-step}
\end{equation}
%
where
$\chi_\Phi
=\operatorname{sgn}(\partial\omega_{01}/\partial\Phi)
\in\{-1,+1\}$
specifies the known local monotonic direction. The applied bias $\Phi_1$ is
restricted to the allowed interval $[\Phi_{\min},\Phi_{\max}]$.

Let $\Delta_{\mathrm{blind}}^{\max}$ denote an experimentally established
upper bound on the frequency displacement produced by the applied blind
step. The proposed initialization is admitted to \textsc{Track} only if
%
\begin{equation}
z_{1-\beta}\sigma_{\Delta,\mathrm{seed}}
+
\Delta_{\mathrm{blind}}^{\max}
\leq
\Delta_{\mathrm{val}},
\label{eq:track-entry-condition}
\end{equation}
%
where $\sigma_{\Delta,\mathrm{seed}}$ includes the uncertainty of the acquired
frequency and the implemented seed reference. This condition reserves enough
detuning margin for the first short-pulse measurement after the blind step.
The determination of $\Delta_{\mathrm{blind}}^{\max}$ and the choice of
$\delta\Phi_{\mathrm{blind}}$ are described in
Appendix~\ref{app:methods}.

No measured sensitivity is available before this step, so the first
short-pulse probe at $\Phi_1$ uses
$\omega_{d,1}=\widehat\omega_{01,\mathrm{seed}}$. Its output supplies the
second frequency estimate $\widehat\omega_{01,1}$ and hence the initial
secant sensitivity
%
\begin{equation}
\widehat S_{\Phi,1}
=
\frac{
\widehat\omega_{01,1}
-
\widehat\omega_{01,\mathrm{seed}}
}{
\Phi_1-\Phi_{\mathrm{seed}}
}.
\label{eq:initial-secant-sensitivity}
\end{equation}
%
This initial slope combines the scanned-Ramsey seed with the short-pulse
estimate at $\Phi_1$. The uncertainty assigned to
$\widehat S_{\Phi,1}$ includes the uncertainties of both frequency estimates.

For subsequent iterations, the local sensitivity is estimated from the two
most recent \textsc{Track} measurements,
%
\begin{equation}
\widehat S_{\Phi,n}
=
\frac{
\widehat\omega_{01,n}
-
\widehat\omega_{01,n-1}
}{
\Phi_n-\Phi_{n-1}
},
\qquad n\geq2.
\label{eq:secant-sensitivity}
\end{equation}
%
This measurement-derived quantity approximates the local
flux-to-frequency sensitivity without supplying the controller with the
quantitative transmon dispersion or its derivative. A sensitivity estimate is
accepted only when its sign is consistent with $\chi_\Phi$ and its magnitude
and uncertainty satisfy the prescribed reliability thresholds. Failure of
this check produces the local-failure event defined in Sec.~II.C.

At each accepted \textsc{Track} iteration, the controller forms
%
\begin{equation}
\widehat r_n
=
\widehat\omega_{01,n}
-
\omega_{\mathrm{tar}}
\label{eq:estimated-feedback-residual}
\end{equation}
%
and tests Eq.~\eqref{eq:target-entry-condition}. If that condition is not
satisfied, the controller proposes the damped secant--Newton update
%
\begin{equation}
\Phi_{n+1}^{\mathrm{prop}}
=
\Phi_n
-
\eta
\frac{\widehat r_n}{\widehat S_{\Phi,n}},
\label{eq:flux-bias-update}
\end{equation}
%
where $\eta\in(0,1]$ is the damping factor. Before application, the commanded
bias change is capped at $\delta\Phi_{\max}$ in magnitude, and the resulting
bias is restricted to $[\Phi_{\min},\Phi_{\max}]$. The uncertainty
propagation, sensitivity thresholds, limiting convention, and controller
parameters are specified in Appendix~\ref{app:methods}.

The flux-bias update changes the physical working point and is intended to
reduce the residual relative to the fixed target. It does not by itself
ensure that the short-pulse probe remains within its validated detuning
interval. The drive-reference and range controller that enforces this second
requirement is constructed next.


\subsection{Drive-reference tracking and range maintenance}
\label{sec:drive-tracking}

After the flux bias is updated, the detuning of the next measurement is
%
\begin{equation}
\Delta_{n+1}
=
\omega_{01}(\Phi_{n+1})
-
\omega_{d,n+1}.
\label{eq:next-step-detuning}
\end{equation}
%
If the microwave reference is held fixed,
$\omega_{d,n+1}=\omega_{d,n}$, this becomes
%
\begin{equation}
\Delta_{n+1}^{\mathrm{fixed}}
=
\omega_{01}(\Phi_{n+1})
-
\omega_{d,n}.
\label{eq:fixed-drive-detuning}
\end{equation}
%
A short-pulse probe referenced to a fixed drive is therefore forced to span
the accumulated frequency displacement as the physical bias moves. The
calibration trajectory can consequently leave the validated interval even
when each flux-bias update moves the qubit toward the target.

During \textsc{Track}, we instead predict the transition frequency at the
updated bias using the measured secant sensitivity and set
%
\begin{equation}
\omega_{d,n+1}
=
\widehat\omega_{01,n}
+
\widehat S_{\Phi,n}
\left(
\Phi_{n+1}-\Phi_n
\right).
\label{eq:drive-tracking}
\end{equation}
%
The right-hand side is the predicted transition frequency at
$\Phi_{n+1}$ based on measurements available through iteration $n$. This
update uses only estimates already produced by the loop and does not evaluate
Eq.~\eqref{eq:transmon-dispersion}. Substitution into
Eq.~\eqref{eq:next-step-detuning} shows that the next measurement detuning is
the one-step frequency-prediction error rather than the accumulated frequency
displacement. Drive tracking therefore reduces the risk of leaving the
validated interval when the local secant prediction remains reliable, but it
does not by itself guarantee range validity.

Let $\sigma_{\mathrm{pred},n+1}$ denote the uncertainty of the frequency
prediction on the right-hand side of
Eq.~\eqref{eq:drive-tracking}. It includes the uncertainties of the current
frequency estimate, the secant sensitivity, and the applied bias step. Before
the next short-pulse probe is executed, the prediction must satisfy
%
\begin{equation}
z_{1-\beta}
\sigma_{\mathrm{pred},n+1}
\leq
\Delta_{\mathrm{val}}.
\label{eq:predicted-range-condition}
\end{equation}
%
For the initial blind step, the corresponding admission test is
Eq.~\eqref{eq:track-entry-condition}. Failure of either premeasurement test
produces a local-failure event and invokes \textsc{Reacquire} before the local
inverse is used.

After an admitted short-pulse measurement, the resulting estimate must also
satisfy
%
\begin{equation}
|\widehat\Delta_{n+1}|
+
z_{1-\beta}\sigma_{\Delta,n+1}
\leq
\Delta_{\mathrm{val}}.
\label{eq:track-range-condition}
\end{equation}
%
This postmeasurement condition checks the consistency of the measured
detuning with the response-validation interval. Failure of this condition,
numerical convergence, inverse-branch continuity, or the sensitivity check
rejects the local estimate and returns the workflow to
\textsc{Reacquire}.

When the range and estimator checks remain satisfied, the resulting frequency
estimate is passed to the next flux-bias update. Satisfaction of
Eq.~\eqref{eq:target-entry-condition} instead produces a target-entry event
and transfers the workflow to \textsc{Lock}, where calibration success is
determined only by the independent verification condition
Eq.~\eqref{eq:verification-condition}.

The flux-bias and drive-reference updates have different physical roles.
Equation~\eqref{eq:flux-bias-update} changes the qubit working point, whereas
Eq.~\eqref{eq:drive-tracking} changes only the reference of the next
measurement. The drive update neither redefines the target nor removes the
calibration residual, because every frequency estimate remains compared with
the fixed $\omega_{\mathrm{tar}}$. Together, the two updates implement the
local \textsc{Track} state defined in Sec.~II.C.

% \subsection{Flux-bias update}
% \label{sec:calibration-protocol}

% The four-state \textsc{Acquire--Track--Lock--Reacquire} workflow and the
% role of each state were defined in Sec.~II.C. An \textsc{Acquire} or \textsc{Reacquire} measurement at
% $\Phi_{\mathrm{seed}}$ supplies the wide-range frequency estimate
% $\widehat\omega_{01,\mathrm{seed}}$ and its uncertainty
% $\sigma_{\omega,\mathrm{seed}}$. This estimate initializes the drive
% reference,
% $\omega_{d,\mathrm{seed}}
% =\widehat\omega_{01,\mathrm{seed}}$.
% The short-pulse estimator is admitted only within the validated local
% interval defined by Eq.~\eqref{eq:validated-interval}. Because the initialized
% drive centers the estimated seed detuning at zero, entry into
% \textsc{Track} requires
% %
% \begin{equation}
% |\widehat\Delta_{\mathrm{seed}}|
% +
% z_{1-\beta}\sigma_{\Delta,\mathrm{seed}}
% \leq
% \Delta_{\mathrm{val}},
% \label{eq:track-entry-condition}
% \end{equation}
% %
% where
% $\widehat\Delta_{\mathrm{seed}}
% =\widehat\omega_{01,\mathrm{seed}}-\omega_{d,\mathrm{seed}}$,
% $\sigma_{\Delta,\mathrm{seed}}$ is its uncertainty, and
% $z_{1-\beta}$ specifies the prescribed confidence margin. This is a
% range-admission condition: it bounds the uncertainty in the initial
% measurement detuning and is distinct from the target-residual condition used
% to enter \textsc{Lock}.

% Within \textsc{Track}, subsequent frequency estimates are supplied by the
% short-pulse estimator. When the in-loop residual first satisfies the
% prescribed target-entry condition, the protocol enters \textsc{Lock}.
% Calibration success is declared only after $N_{\mathrm{lock}}$ consecutive
% independent verification measurements confirm the target condition. A failed
% confirmation returns the protocol to \textsc{Track} if the local validity
% conditions remain satisfied; otherwise it triggers \textsc{Reacquire}. A
% range violation, a discontinuous inverse, insufficient confidence, an
% anomalous local sensitivity, or repeated measurement failure during
% \textsc{Track} also triggers \textsc{Reacquire}. The resulting wide-range
% estimate may then seed a new \textsc{Track} episode.

% The accepted \textsc{Acquire} or \textsc{Reacquire} result is reused as the
% initial measured pair
% $(\Phi_{\mathrm{seed}},\widehat\omega_{01,\mathrm{seed}})$; no additional
% short-pulse measurement is performed at the seed bias. The controller forms
% the corresponding target residual
% %
% \begin{equation}
% \widehat r_{\mathrm{seed}}
% =
% \widehat\omega_{01,\mathrm{seed}}
% -
% \omega_{\mathrm{tar}}.
% \label{eq:seed-residual}
% \end{equation}
% %
% If this residual already satisfies the target-entry condition, the protocol
% proceeds directly to \textsc{Lock} for independent confirmation. Otherwise,
% the controller proposes a second bias,
% %
% \begin{equation}
% \Phi_1^{\mathrm{prop}}
% =
% \Phi_{\mathrm{seed}}
% -
% \chi_\Phi
% \operatorname{sgn}(\widehat r_{\mathrm{seed}})
% \delta\Phi_{\mathrm{blind}},
% \label{eq:blind-step}
% \end{equation}
% %
% where
% $\chi_\Phi=\operatorname{sgn}(\partial\omega_{01}/\partial\Phi)
% \in\{-1,+1\}$
% specifies the known local monotonic direction. The applied value $\Phi_1$ is
% restricted to the allowed interval $[\Phi_{\min},\Phi_{\max}]$. This
% initialization uses one acquired seed followed by a bounded one-sided blind
% step and assumes neither a quantitative local sensitivity nor a symmetric
% two-point bracket.

% Because no measured sensitivity is yet available, the first short-pulse
% probe at $\Phi_1$ uses the acquired frequency as its drive reference,
% $\omega_{d,1}=\widehat\omega_{01,\mathrm{seed}}$. The resulting estimate
% $\widehat\omega_{01,1}$ supplies the second bias--frequency pair. The initial
% secant sensitivity is therefore
% %
% \begin{equation}
% \widehat S_{\Phi,1}
% =
% \frac{
% \widehat\omega_{01,1}
% -
% \widehat\omega_{01,\mathrm{seed}}
% }{
% \Phi_1-\Phi_{\mathrm{seed}}
% }.
% \label{eq:initial-secant-sensitivity}
% \end{equation}
% %
% This first estimate combines the scanned-Ramsey seed with the short-pulse
% estimate at $\Phi_1$. Both estimates refer to the same transition frequency;
% their uncertainties, together with any calibrated relative offset between
% the two probes, are included when assessing the reliability of
% $\widehat S_{\Phi,1}$.

% After this initialization, the local sensitivity is estimated from the two
% most recent \textsc{Track} measurements,
% %
% \begin{equation}
% \widehat S_{\Phi,n}
% =
% \frac{
% \widehat\omega_{01,n}
% -
% \widehat\omega_{01,n-1}
% }{
% \Phi_n-\Phi_{n-1}
% },
% \qquad n\geq 2.
% \label{eq:secant-sensitivity}
% \end{equation}
% %
% This measurement-derived quantity approximates the local
% flux-to-frequency sensitivity without supplying the controller with the
% quantitative dispersion or its derivative. At each \textsc{Track} iteration,
% the controller forms
% %
% \begin{equation}
% \widehat r_n
% =
% \widehat\omega_{01,n}
% -
% \omega_{\mathrm{tar}}
% \label{eq:estimated-feedback-residual}
% \end{equation}
% %
% and proposes the damped secant--Newton update
% %
% \begin{equation}
% \Phi_{n+1}^{\mathrm{prop}}
% =
% \Phi_n
% -
% \eta
% \frac{\widehat r_n}{\widehat S_{\Phi,n}},
% \label{eq:flux-bias-update}
% \end{equation}
% %
% where $\eta\in(0,1]$ is the damping factor. Before the proposal is applied,
% the commanded bias change is capped at $\delta\Phi_{\max}$ in magnitude, and
% the resulting bias is restricted to
% $[\Phi_{\min},\Phi_{\max}]$. These safeguards are applied at every
% \textsc{Track} iteration. The explicit limiting convention, uncertainty
% propagation, and controller parameters are given in
% Appendix~\ref{app:methods}.

% The flux-bias update moves the physical working point and is intended to
% reduce the residual relative to the fixed target. It does not, however,
% ensure that the next short-pulse measurement remains within the validated
% detuning interval. The following subsection addresses this second
% requirement by updating the drive reference used for the next probe.

% \subsection{Drive-reference tracking}
% \label{sec:drive-tracking}

% At the next iteration, the measurement detuning is generally
% %
% \begin{equation}
% \Delta_{n+1}
% =
% \omega_{01}(\Phi_{n+1})-\omega_{d,n+1}.
% \label{eq:next-step-detuning}
% \end{equation}
% %
% If the microwave reference is held fixed,
% $\omega_{d,n+1}=\omega_{d,n}$, Eq.~\eqref{eq:next-step-detuning} reduces to
% %
% \begin{equation}
% \Delta_{n+1}^{\mathrm{fixed}}
% =
% \omega_{01}(\Phi_{n+1})-\omega_{d,n}.
% \label{eq:fixed-drive-detuning}
% \end{equation}
% %
% A local probe referenced to a fixed drive is therefore forced to span the
% accumulated frequency displacement as the flux bias moves, and the calibration
% trajectory can leave the validated interval even when the physical working
% point moves toward the target.

% During \textsc{Track}, we instead predict the transition frequency at the
% updated bias using the measured secant sensitivity and set
% %
% \begin{equation}
% \omega_{d,n+1}
% =
% \widehat\omega_{01,n}
% +
% \widehat S_{\Phi,n}
% \left(
% \Phi_{n+1}-\Phi_n
% \right).
% \label{eq:drive-tracking}
% \end{equation}
% %
% This update uses only estimates already produced by the loop and does not
% evaluate Eq.~\eqref{eq:transmon-dispersion}. Substitution into
% Eq.~\eqref{eq:next-step-detuning} shows that the next measurement detuning is
% the one-step frequency-prediction error rather than the accumulated frequency
% displacement. The update therefore reduces the risk of leaving the validated
% interval when the local secant prediction is accurate, but it does not
% guarantee range validity when that prediction fails. It changes only the
% measurement reference: the physical calibration residual remains defined
% relative to the fixed target $\omega_{\mathrm{tar}}$.

% Across the full workflow, \textsc{Acquire} and \textsc{Reacquire} establish
% a wide-range frequency reference, while \textsc{Track} uses
% Eq.~\eqref{eq:drive-tracking} to maintain the local measurement range. A
% failure of the reference prediction, confidence condition, or local inverse
% rejects the local update and returns the protocol to \textsc{Reacquire}. In
% \textsc{Lock}, the drive is placed at $\omega_{\mathrm{tar}}$ while independent
% measurements confirm calibration. Failed verification returns the protocol to
% \textsc{Track} when the local conditions remain valid and to
% \textsc{Reacquire} otherwise.

% \section{Response-kernel calibration}
% \label{sec:kernels}

% Let $U_0(t)$ be the propagator of the control Hamiltonian at zero detuning and
% define the interaction-picture perturbation operator
% \begin{equation}
%  W(t)=\frac{1}{2}U_0^\dagger(t)\sigma_zU_0(t).
% \end{equation}
% For either analysis branch, the final population is a functional of a
% time-dependent detuning.  Its Volterra expansion contains the kernels
% \begin{multline}
%  k_n(t_1,\ldots,t_n)=i^n\bra{0}[W(t_{(n)}),\\
%  [\ldots,[W(t_{(1)}),M_I(T)]\ldots]]\ket{0},
%  \label{eq:nested-kernel}
% \end{multline}
% where $t_{(1)}>\cdots>t_{(n)}$ and $M_I(T)$ is the final measurement operator
% in the control interaction picture.  We form differential kernels from the
% two branches in Eq.~\eqref{eq:sequences}.

% For a constant detuning applied throughout the sequence, the coefficients in
% Eq.~\eqref{eq:cubic-response} are integrals over the complete kernels,
% \begin{align}
%  A_1 &= \int_0^T k_1(t)\,\dd t,\\
%  A_3 &= \int_0^T\!\dd t_1\int_0^T\!\dd t_2\int_0^T\!\dd t_3\,
%  k_3(t_1,t_2,t_3).
%  \label{eq:kernel-integrals}
% \end{align}
% The triple integral includes unequal-time correlations generated by coherent
% evolution.  Replacing it by $\int k_3(t,t,t)\dd t$ discards these off-diagonal
% terms and does not give the cubic coefficient of a constant detuning.

% We use the full three-time kernel for the drive-tracked protocol.  An
% alternative calibration fits an odd polynomial to $p_{\mathrm d}(\Delta)$ over
% a symmetric detuning scan.  The fit is useful as an independent check when the
% drive is fixed, but it assumes that the scan is centered on the symmetry point.
% Drive tracking moves $\omega_d$ at every iteration; a fit performed around a
% different reference can then return an incorrect or even sign-changing linear
% coefficient.  In contrast, Eq.~\eqref{eq:kernel-integrals} is evaluated for the
% actual pulse and drive frequency and does not require a symmetric detuning
% scan.
